\documentclass[aspectratio=169]{beamer}
\usepackage[english]{babel}
\usepackage[utf8]{inputenc}
\usepackage[T1]{fontenc}
\usepackage{lmodern}
\usepackage{listings}
\usepackage{xcolor}
\usepackage{graphicx}
\usepackage{textcomp}
\DeclareUnicodeCharacter{2014}{---}
\DeclareUnicodeCharacter{2013}{--}
\DeclareUnicodeCharacter{2212}{-}
\DeclareUnicodeCharacter{2192}{\ensuremath{\rightarrow}}
\DeclareUnicodeCharacter{00B1}{\ensuremath{\pm}}
\DeclareUnicodeCharacter{00D7}{\ensuremath{\times}}
\DeclareUnicodeCharacter{00B7}{\ensuremath{\cdot}}
\DeclareUnicodeCharacter{20AC}{EUR}
\DeclareUnicodeCharacter{2070}{\textsuperscript{0}}
\DeclareUnicodeCharacter{00B9}{\textsuperscript{1}}
\DeclareUnicodeCharacter{00B2}{\textsuperscript{2}}
\DeclareUnicodeCharacter{00B3}{\textsuperscript{3}}
\DeclareUnicodeCharacter{2074}{\textsuperscript{4}}
\DeclareUnicodeCharacter{2075}{\textsuperscript{5}}
\DeclareUnicodeCharacter{2076}{\textsuperscript{6}}
\DeclareUnicodeCharacter{2077}{\textsuperscript{7}}
\DeclareUnicodeCharacter{2078}{\textsuperscript{8}}
\DeclareUnicodeCharacter{2079}{\textsuperscript{9}}
\DeclareUnicodeCharacter{2248}{\ensuremath{\approx}}
\DeclareUnicodeCharacter{2264}{\ensuremath{\le}}

% Colors for code
\definecolor{codebg}{RGB}{245,245,245}
\definecolor{codecomment}{RGB}{0,128,0}
\definecolor{codekeyword}{RGB}{0,0,255}
\definecolor{codestring}{RGB}{163,21,21}
\definecolor{bandcolor}{RGB}{200,50,50}

\lstdefinestyle{pythonstyle}{
    language=Python,
    backgroundcolor=\color{codebg},
    basicstyle=\ttfamily\scriptsize,
    keywordstyle=\color{codekeyword},
    commentstyle=\color{codecomment},
    stringstyle=\color{codestring},
    showstringspaces=false,
    breaklines=true,
    frame=single,
    rulecolor=\color{black},
    escapeinside={(*@}{@*)},
    moredelim=**[l][\color{bandcolor}\bfseries]{\#\ ---\ NEW},
    moredelim=**[l][\color{bandcolor}\bfseries]{\#\ ---\ CHANGED},
    literate=
      {á}{{\'a}}1 {é}{{\'e}}1 {í}{{\'i}}1 {ó}{{\'o}}1 {ú}{{\'u}}1
      {Á}{{\'A}}1 {É}{{\'E}}1 {Í}{{\'I}}1 {Ó}{{\'O}}1 {Ú}{{\'U}}1
      {ñ}{{\~n}}1 {Ñ}{{\~N}}1 {ü}{{\"u}}1 {Ü}{{\"U}}1
      {¿}{{?`}}1 {¡}{{!`}}1
      {—}{{---}}1 {–}{{--}}1 {−}{{-}}1
      {±}{{\ensuremath{\pm}}}1 {×}{{\ensuremath{\times}}}1
      {→}{{\ensuremath{\rightarrow}}}1
      {°}{{\textdegree}}1
      {·}{{\ensuremath{\cdot}}}1
      {€}{{\texteuro}}1
      {≈}{{\ensuremath{\approx}}}1
      {≤}{{\ensuremath{\le}}}1
      {⁰}{{\textsuperscript{0}}}1 {¹}{{\textsuperscript{1}}}1
      {²}{{\textsuperscript{2}}}1 {³}{{\textsuperscript{3}}}1
      {⁴}{{\textsuperscript{4}}}1 {⁵}{{\textsuperscript{5}}}1
      {⁶}{{\textsuperscript{6}}}1 {⁷}{{\textsuperscript{7}}}1
      {⁸}{{\textsuperscript{8}}}1 {⁹}{{\textsuperscript{9}}}1
}

\lstset{style=pythonstyle}

% Title slide macro: \lessontitle{Lesson Number}{Title}
\newcommand{\lessontitle}[2]{
    \title{Lesson #1: #2}
    \author{}
    \institute{Tec de Monterrey}
    \begin{frame}[plain]
        \titlepage
    \end{frame}
}

% Figure frame macro: \figureframe{Title}{Image path}
\newcommand{\figureframe}[2]{
    \begin{frame}{#1}
        \begin{center}
            \includegraphics[height=0.8\textheight,width=\textwidth,keepaspectratio]{#2}
        \end{center}
    \end{frame}
}


% Listings pulled from src/ with \lstinputlisting, so the slide is the file.
\lstset{
    basicstyle=\ttfamily\fontsize{8pt}{9.6pt}\selectfont,
    breaklines=true,
    breakatwhitespace=true,
    columns=fullflexible,
    keepspaces=true,
    xleftmargin=0.3em,
    framesep=2pt,
    aboveskip=0pt,
    belowskip=0pt
}

\newcommand{\resultframe}[3]{%
    \begin{frame}{#1}
        \begin{center}
            \includegraphics[height=0.62\textheight,width=\textwidth,keepaspectratio]{#2}
        \end{center}
        \vspace{0.25em}
        {\small #3\par}
    \end{frame}%
}

\newcommand{\claimframe}[2]{%
    \begin{frame}{#1}
        {\small #2\par}
    \end{frame}%
}

\date{}

\begin{document}

\lessontitle{11}{Other Problems}

\section{This lesson}

\begin{frame}{This lesson}
Lesson 08 established that you usually cannot tell whether a run succeeded. These two problems are the rare case where one check settles it.

\vspace{0.6em}
\begin{itemize}
    \item A verifiable objective can be a stopping rule (residual 0, conflicts 0).
    \item When the constraint is the objective, the landscape near the answer is a plateau.
    \item A finite integer box can be enumerated; that answers questions a GA cannot.
\end{itemize}
\end{frame}

% ================= colouring_01_the_graph =================
\begin{frame}{Witnesses and certificates}
\[
R(x,y,z)=|f|+|g|+|w|,\quad (x,y,z)\in[-20,20]^3\cap\mathbb Z^3,
\quad |\mathcal X|=41^3=68{,}921,
\]
\[
C(c)=\sum_{(u,v)\in E}I[c_u=c_v],\quad
|\mathcal C|=3^{30}=205{,}891{,}132{,}094{,}649,\quad
\mathbb E[C]=|E|/3.
\]
Zero residual or conflicts is a witness. Enumeration and exact backtracking add
certificates of uniqueness or impossibility that a successful GA run cannot.
\end{frame}

\section{When the constraint is the whole objective}

\begin{frame}{When the constraint is the whole objective}
Graph colouring is the opposite of the equation system. There the objective was
a smooth-ish quantity to drive down; here the only thing being measured is a
violated constraint. Fitness is the negated number of edges whose two ends share
a colour, so a legal colouring and an optimal colouring are the same object and
there is nothing to optimise once you reach it.
\end{frame}

\resultframe{When the constraint is the whole objective}{../figures/colouring_01_the_graph.png}{%
    30 vertices, 64 distinct edges, mean degree 4.27, and 3³⁰ = 205,891,132,094,649 colourings. A random colouring at seed 7 breaks 26 of the 64 edges}

% ================= colouring_02_random_population =================
\section{The baseline luck provides}

\begin{frame}{The baseline luck provides}
A random colouring assigns each of the three colours independently, so each edge
has one chance in three of being broken. That gives an expected conflict count
of |E|/3 before running anything, and the sample should land on it. What matters
is the other end of the histogram: the best of 500 random colourings, and how
far it still is from the 0 that the problem demands.
\end{frame}

\begin{frame}[fragile]{(1) random\_population()}
\lstinputlisting[basicstyle=\ttfamily\fontsize{8pt}{9.6pt}\selectfont,firstline=86,lastline=90]{../src/colouring_02_random_population.py}
\end{frame}

\begin{frame}[fragile]{(2) the sample}
\lstinputlisting[basicstyle=\ttfamily\fontsize{8pt}{9.6pt}\selectfont,firstline=93,lastline=101]{../src/colouring_02_random_population.py}
\end{frame}

\resultframe{The baseline luck provides}{../figures/colouring_02_random_population.png}{%
    /3 = 21.33`, and the best draw still breaks 11 edges (17.2\%)}

% ================= colouring_03_fitness_driven_operators =================
\section{Operators that consult the fitness function}

\begin{frame}{Operators that consult the fitness function}
On a plateau, a blind operator is a coin toss. Gridin's answer for this problem
is to let the operators look before they leap: a crossover that keeps the best
two of \{both children, both parents\}, and a mutation that tries a few random
recolourings and keeps one only if it lowers the conflict count. Both are
greedy, and greed is a real choice with a real cost, so measure it rather than
assert it.
\end{frame}

\begin{frame}[fragile]{(1) crossover\_n\_point()}
\lstinputlisting[basicstyle=\ttfamily\fontsize{8pt}{9.6pt}\selectfont,firstline=95,lastline=107]{../src/colouring_03_fitness_driven_operators.py}
\end{frame}

\begin{frame}[fragile]{(2) crossover\_fitness\_driven\_n\_point()}
\lstinputlisting[basicstyle=\ttfamily\fontsize{8pt}{9.6pt}\selectfont,firstline=110,lastline=118]{../src/colouring_03_fitness_driven_operators.py}
\end{frame}

\begin{frame}[fragile]{(3) mutation\_fitness\_driven\_random\_change()}
\lstinputlisting[basicstyle=\ttfamily\fontsize{8pt}{9.6pt}\selectfont,firstline=121,lastline=133]{../src/colouring_03_fitness_driven_operators.py}
\end{frame}

\begin{frame}[fragile]{(4) the experiment}
\lstinputlisting[basicstyle=\ttfamily\fontsize{7pt}{8.5pt}\selectfont,firstline=136,lastline=162]{../src/colouring_03_fitness_driven_operators.py}
\end{frame}

\resultframe{Operators that consult the fitness function}{../figures/colouring_03_fitness_driven_operators.png}{%
    Plain two-point crossover moves the population mean by +0.03 edges; the fitness-driven version by −2.32. Greedy recolouring moves the mean by −1.02 but leaves 206 of 500 individuals untouched, and its best individual (11) is *worse* than a blind recolour's best (9)}

% ================= colouring_04_the_full_search =================
\section{The search that does not finish, and the one that does}

\begin{frame}{The search that does not finish, and the one that does}
This is the step the lesson is built for. The genetic algorithm is assembled and
run, and on most seeds it stops one or two edges short of a legal colouring and
stays there for the rest of its budget - the plateau that step 1 predicted.
Then ten lines of backtracking settle the same question outright, and also prove
that two colours are impossible, which no amount of running the GA could have
established.
\end{frame}

\begin{frame}[fragile]{(1) selection\_rank\_with\_elite()}
\lstinputlisting[basicstyle=\ttfamily\fontsize{7pt}{8.5pt}\selectfont,firstline=139,lastline=158]{../src/colouring_04_the_full_search.py}
\end{frame}

\begin{frame}[fragile]{(2) run()}
\lstinputlisting[basicstyle=\ttfamily\fontsize{7pt}{8.5pt}\selectfont,firstline=161,lastline=187]{../src/colouring_04_the_full_search.py}
\end{frame}

\begin{frame}[fragile]{(3) exact\_colouring()}
\lstinputlisting[basicstyle=\ttfamily\fontsize{6pt}{7.2pt}\selectfont,firstline=190,lastline=227]{../src/colouring_04_the_full_search.py}
\end{frame}

\begin{frame}[fragile]{(4) the seed sweep}
\lstinputlisting[basicstyle=\ttfamily\fontsize{8pt}{9.6pt}\selectfont,firstline=230,lastline=241]{../src/colouring_04_the_full_search.py}
\end{frame}

\resultframe{The search that does not finish, and the one that does}{../figures/colouring_04_the_full_search.png}{%
    \textbf{1 of 4 seeds} reaches a legal colouring (seed 11, generation 15). Seeds 7 and 16 stall at 2 conflicts and seed 1 at 1 conflict, all still stuck after the full 200 generations. Backtracking finds a proper 3-colouring in 40 search nodes and proves no 2-colouring exists in 4 nodes — 0.69 ms for both}

% ================= equations_01_the_system =================
\section{A problem whose answer can be checked}



\resultframe{A problem whose answer can be checked}{../figures/equations_01_the_system.png}{%
    Genes are integers in \texttt{[-20, 20]}, so the whole space is 68,921 triples. Residuals are integers, not floats: along the line \texttt{(y, z) = (1, 1)} they run from 3 digits at \texttt{x = 2} to 31 digits at \texttt{x = 20}}

% ================= equations_02_random_population =================
\section{What a random sample is worth here}

\begin{frame}{What a random sample is worth here}
Before running a genetic algorithm, find out what luck alone buys. In lesson 09
random sampling was a respectable baseline. Here it is not, and the reason is
visible in one histogram: the residual is measured in decimal digits, so a draw
that looks close is still astronomically far away.
\end{frame}

\begin{frame}[fragile]{(1) random\_triple()}
\lstinputlisting[basicstyle=\ttfamily\fontsize{8pt}{9.6pt}\selectfont,firstline=75,lastline=79]{../src/equations_02_random_population.py}
\end{frame}

\begin{frame}[fragile]{(2) the sample}
\lstinputlisting[basicstyle=\ttfamily\fontsize{8pt}{9.6pt}\selectfont,firstline=82,lastline=92]{../src/equations_02_random_population.py}
\end{frame}

\resultframe{What a random sample is worth here}{../figures/equations_02_random_population.png}{%
    400 uniform draws (0.58\% of the box) find 0 solutions; the best is \texttt{(0, 1, 2)} with residual 453, while the median draw has a 957-digit residual and the worst 14,187 digits}

% ================= equations_03_the_full_ga =================
\section{The genetic algorithm solves it}

\begin{frame}{The genetic algorithm solves it}
Nothing here is new machinery: rank selection with elitism came from lesson 03,
blend crossover from lesson 04, Gaussian deviation from lesson 05. What is new
is the stopping rule. Every other lesson stopped after a fixed budget because it
could not tell whether it had won. This one stops when the residual is zero,
because that is a proof.
\end{frame}

\begin{frame}[fragile]{(1) Individual}
\lstinputlisting[basicstyle=\ttfamily\fontsize{8pt}{9.6pt}\selectfont,firstline=88,lastline=99]{../src/equations_03_the_full_ga.py}
\end{frame}

\begin{frame}[fragile]{(2) selection\_rank\_with\_elite()}
\lstinputlisting[basicstyle=\ttfamily\fontsize{7pt}{8.5pt}\selectfont,firstline=102,lastline=123]{../src/equations_03_the_full_ga.py}
\end{frame}

\begin{frame}[fragile]{(3) crossover\_blend()}
\lstinputlisting[basicstyle=\ttfamily\fontsize{8pt}{9.6pt}\selectfont,firstline=126,lastline=139]{../src/equations_03_the_full_ga.py}
\end{frame}

\begin{frame}[fragile]{(4) mutation\_random\_deviation()}
\lstinputlisting[basicstyle=\ttfamily\fontsize{8pt}{9.6pt}\selectfont,firstline=142,lastline=151]{../src/equations_03_the_full_ga.py}
\end{frame}

\begin{frame}[fragile]{(5) run()}
\lstinputlisting[basicstyle=\ttfamily\fontsize{6pt}{7.2pt}\selectfont,firstline=154,lastline=187]{../src/equations_03_the_full_ga.py}
\end{frame}

\resultframe{The genetic algorithm solves it}{../figures/equations_03_the_full_ga.png}{%
    Seed 3 reaches residual 0 at generation 5 after 2,794 evaluations, at \texttt{x = -6, y = 2, z = 3}; \texttt{f}, \texttt{g} and \texttt{w} are each individually 0}

% ================= equations_04_exhaustive_search =================
\section{The method that does not need a genetic algorithm}

\begin{frame}{The method that does not need a genetic algorithm}
The genes are integers and the box is finite, so the entire search space can be
walked. This is the step where the course finally asks the question it has been
avoiding: given that a genetic algorithm is a heuristic, when is it the wrong
tool? The answer here is not the one the evaluation counts suggest, and the
step is written around what the code actually measures.
\end{frame}

\begin{frame}[fragile]{(1) enumerate\_box()}
\lstinputlisting[basicstyle=\ttfamily\fontsize{7pt}{8.5pt}\selectfont,firstline=179,lastline=199]{../src/equations_04_exhaustive_search.py}
\end{frame}

\begin{frame}[fragile]{(2) run()}
\lstinputlisting[basicstyle=\ttfamily\fontsize{8pt}{9.6pt}\selectfont,firstline=145,lastline=145]{../src/equations_04_exhaustive_search.py}
\end{frame}

\begin{frame}[fragile]{(3) the comparison}
\lstinputlisting[basicstyle=\ttfamily\fontsize{8pt}{9.6pt}\selectfont,firstline=202,lastline=216]{../src/equations_04_exhaustive_search.py}
\end{frame}

\resultframe{The method that does not need a genetic algorithm}{../figures/equations_04_exhaustive_search.png}{%
    All 5 seeds tried find the same root, in 5–28 generations and 2,794–13,764 evaluations. Enumerating the box costs 68,921 evaluations (about 10× more, roughly 3.7 s) and returns strictly more: there is \textbf{exactly one} root in the box}

\section{Next}

\begin{frame}{Next}
The equation system has exactly one root in the box. The GA finds it cheaper than enumeration --- and still cannot prove uniqueness. Colouring has 205 trillion assignments, yet exact backtracking certifies a 3-colouring and rules out 2 colours.

\vspace{0.8em}
\textbf{Lesson 12 --- The Adaptive Genetic Algorithm}

Lesson 07's tuned fixed settings, Lesson 10's TSP, equal evaluation budget. Adaptation is insurance against a bad constant, not a replacement for tuning.
\end{frame}
\end{document}
